\section{Conclusion}

We followed the applicability conditions of answers through a chunked RAG pipeline on ConditionalQA, under an analysis plan fixed before any evaluation run. Conditions are lost at every stage, but not equally. Chunking isolates the conditions of many questions in chunks of their own, and losing such a chunk costs every generator we tested, under every prompt we tried. Retrieval finds conditions as often as answers, costs more through the distractors it adds than through what it misses, and four context-preserving retrieval designs pass on no more conditions at the same budget. The evaluated evidence-reducing operators remove condition chunks as readily as answer chunks, since nothing in their design tells an exception from a contradiction or an irrelevant passage, and on RAMDocs they do not single out misinformation either. Generation shows the largest observed deficit: with every condition in context, five generators from 7B to 550B parameters state at most a third of the conditions, a figure that lenient or meaning-based matching barely changes, and six prompting strategies, including chain-of-condition reasoning, move it only within a range that leaves most conditions unstated, while those that state more attach conditions to unconditional answers too. The stages are measured on different scales, so this ranks the observed deficits rather than decomposing the end-to-end loss. Our qualifier-aware completion gains at most about a point and a half under either threshold protocol, not significantly in the confirmatory test, and even a perfect scorer would gain at most a few points.

Three consequences follow for practice. First, evidence-reducing and conflict-handling modules should be evaluated on data where the evidence they remove can be a valid condition, not only where it is false, and with strong generators, which lose at least as much before generation. Second, retrieval and chunking choices are not where most conditions are lost; keeping conditions with their rules helps only as far as the generator states what it is given. Third, since the official metric favours a generator that states no conditions, scores should be reported separately for conditional and unconditional questions.

\textbf{Limitations.} The benchmark's human agreement is 56.6 F1 with conditions, so part of what we count as loss is annotation noise; no human-validated sample of the metric's false negatives was drawn, although token-overlap and NLI-based matching both leave the conclusion unchanged. The two primary generators are small and 4-bit quantised; the three larger ones answer a stratified subsample, were served by more than one provider, and include two models chosen after the plan was written, recorded as deviations before their answers were analysed. One prompt was frozen for every system; six strategies do not span every reasonable prompt, so the generation finding is a claim about the strategies evaluated. Each operator is our implementation of a published method's evidence-reducing step, not the complete system. All texts are English UK public-service guidance, a genre of rules, lists and exceptions; without a second annotated domain, our quantitative results are claims about structured policy guidance, and healthcare, financial and legal documents are natural extensions. RAMDocs mixes conflicting and ambiguous questions, and its misinformation is synthetic; the transfer corpus is small and was labelled by the authors.

The clearest direction for future work is condition-aware generation: decoding or training objectives that make a reader state, rather than ignore, the conditions already in its context, and that teach it to judge whether the scenario already settles them; our verification pass and prompt variants show that the first without the second does more harm than good.
